\chapter{Teeth keeping them clean}
\index{Teeth keeping them clean}

\medskip
\noindent\textit{Educational reference; always seek professional advice for individual care.}

\section*{Definition and Overview}
Keeping teeth clean means removing dental plaque---the sticky biofilm of bacteria that forms continuously on tooth surfaces---before it hardens into calculus (tartar) or produces the acids and toxins that cause tooth decay and gum disease. Effective cleaning is the foundation of lifelong oral health: it protects enamel, keeps gums firm and pink, reduces bad breath, and lowers the risk of tooth loss. Cleaning is not only brushing; it includes interdental cleaning (floss, brushes, or water flossers), fluoride use, sensible diet, and professional care. In many countries, public health messaging from the local Dental Association and state health services emphasises twice-daily fluoride toothpaste, limiting free sugars, and regular dental visits. This chapter explains practical, evidence-based methods suitable for children and adults, common mistakes, and when professional help is needed.

\section*{Epidemiology}
Dental caries and periodontal disease remain among the most prevalent chronic conditions worldwide. In many countries, a substantial proportion of children still experience decay in primary or permanent teeth, with higher rates in some rural, remote, and socioeconomically disadvantaged communities and among Aboriginal and Torres Strait Islander peoples. Adults commonly have untreated decay, fillings, or gum inflammation; severe periodontitis affects a meaningful minority and is a leading cause of tooth loss. Fluoridated water supplies cover much of the population and have reduced decay at a community level, but individual plaque control still determines much of personal risk. People with dry mouth (from medicines, radiotherapy, or Sjögren syndrome), diabetes, smoking, orthodontic appliances, or disability-related care barriers need extra support to keep teeth clean.

\section*{Aetiology and Pathophysiology}
Plaque bacteria metabolise dietary sugars and starches, producing acids that demineralise enamel (caries) and inflammatory products that trigger gingivitis and, over time, periodontitis with bone loss.

\begin{itemize}
    \item \textbf{Biofilm formation:} salivary proteins coat the tooth within minutes; bacteria attach and mature into organised plaque within hours if not disrupted
    \item \textbf{Acid attack and remineralisation balance:} frequent sugar exposure lowers pH; saliva and fluoride help repair early mineral loss if attacks are limited
    \item \textbf{Gingival inflammation:} plaque at the gum margin causes redness, swelling, and bleeding (gingivitis), which is reversible with cleaning
    \item \textbf{Periodontal pocketing:} in susceptible people, chronic inflammation destroys attachment and bone, creating pockets that trap more plaque
    \item \textbf{Calculus:} mineralised plaque provides a rough surface that holds fresh biofilm and cannot be removed by brushing alone
    \item \textbf{Dry mouth and appliances:} reduced saliva and brackets, dentures, or implants change cleaning needs and raise caries or candida risk
\end{itemize}

\section*{Clinical Features}
\begin{itemize}
    \item \textbf{Healthy cleaned mouth:} firm pink gums that do not bleed easily, clean tooth surfaces, fresh breath, and comfortable chewing
    \item \textbf{Early neglect signs:} bleeding on brushing or flossing, bad taste, fuzzy tooth feel, and mild gum swelling
    \item \textbf{Caries-related clues:} sensitivity to sweet or cold, visible white or brown spots, food packing, or holes
    \item \textbf{Periodontal clues:} receding gums, longer-looking teeth, mobility, pus, or drifting teeth
    \item \textbf{Halitosis:} persistent odour often from tongue coating, interdental plaque, or gum pockets
    \item \textbf{Staining:} coffee, tea, red wine, and tobacco stain plaque and rough surfaces; professional cleaning removes much extrinsic stain
\end{itemize}

Bleeding gums are a sign to clean more effectively, not to stop cleaning---avoidance worsens inflammation.

\section*{Differential Diagnosis}
\begin{itemize}
    \item \textbf{Gingivitis versus periodontitis:} bleeding alone may be gingivitis; true attachment loss needs professional probing and radiographs
    \item \textbf{Medication-related gum overgrowth:} calcium channel blockers, phenytoin, and ciclosporin can enlarge gums and trap plaque
    \item \textbf{Hormonal gingivitis:} pregnancy can exaggerate inflammatory response to the same plaque load
    \item \textbf{Vitamin C deficiency and systemic disease:} uncommon but relevant in restricted diets or severe illness
    \item \textbf{Oral thrush or viral ulcers:} white patches or sores that are not removed by brushing need different care
    \item \textbf{Dry mouth caries pattern:} root and cervical decay despite apparently reasonable brushing
\end{itemize}

\section*{When to Seek Medical Care}
Book a dentist for routine checks and cleans as advised (often six to twelve monthly, individualised), sooner for pain, swelling, trauma, or bleeding that does not improve with a fortnight of careful home care. See a GP for uncontrolled diabetes, suspected eating disorders, or medicines causing severe dry mouth.

\textbf{Contact your local emergency services for:}
\begin{itemize}
    \item Facial swelling affecting the eye or breathing, or rapidly spreading infection with fever
    \item Uncontrolled oral bleeding or trauma with suspected jaw fracture
    \item Severe allergic reaction after dental products or antibiotics
\end{itemize}

\section*{Diagnosis and Investigations}
\begin{itemize}
    \item \textbf{Clinical dental exam:} plaque and calculus scores, gum probing depths, bleeding points, mobility, and caries detection
    \item \textbf{Radiographs:} bitewings and periapicals to find hidden decay and bone levels
    \item \textbf{Oral hygiene review:} demonstration of brushing and interdental technique; disclosing tablets to show missed plaque
    \item \textbf{Saliva and risk assessment:} flow, diet diary, fluoride exposure, smoking, and medical history
    \item \textbf{Periodontal charting:} baseline and monitoring for gum disease
\end{itemize}

\section*{Management}
Daily self-care plus professional support keeps most mouths healthy.

\begin{itemize}
    \item \textbf{Brushing:} twice daily for two minutes with a soft toothbrush (manual or powered) and fluoride toothpaste (at least 1000 ppm fluoride for children as age-directed; 1450 ppm commonly for adults in many countries); spit, do not rinse heavily, to leave fluoride on teeth
    \item \textbf{Interdental cleaning:} floss, flossettes, or interdental brushes sized to the gaps, once daily; water flossers help some people with braces or dexterity limits
    \item \textbf{Tongue cleaning:} gentle tongue scraping or brushing reduces coating and odour
    \item \textbf{Fluoride extras:} high-fluoride toothpaste or professional varnish for high caries risk, as prescribed
    \item \textbf{Professional scale and clean:} removes calculus and stain; root surface debridement for periodontitis
    \item \textbf{Diet advice:} fewer sugar hits per day; water between meals; limit sticky sweets and sugary drinks including fruit juice
    \item \textbf{Smoking cessation:} critical for gum health and cancer risk reduction
    \item \textbf{Denture and implant care:} remove and clean dentures daily; clean around implants with appropriate brushes without damaging tissues
    \item \textbf{Children:} parental brushing until technique is reliable (often around eight years); pea-sized fluoride toothpaste; supervised spitting
\end{itemize}

\section*{Complications}
\begin{itemize}
    \item \textbf{Dental caries:} pain, infection, abscess, and tooth loss if untreated
    \item \textbf{Periodontitis:} bone loss, mobility, abscess, and systemic inflammatory associations including poorer diabetes control
    \item \textbf{Tooth loss and impaired nutrition:} difficulty chewing healthy foods
    \item \textbf{Cost and quality-of-life burden:} repeated restorations, dentures, and social impact of appearance and breath
    \item \textbf{Acute dental infections:} can rarely spread to deep neck spaces and become life-threatening
\end{itemize}

\section*{Prognosis and Outcomes}
Gingivitis reverses within days to weeks of effective plaque removal. Early caries can remineralise with fluoride and diet change. Established cavities need restoration. Periodontitis can usually be stabilised with professional care and excellent home cleaning, though lost bone does not fully regrow without specialised procedures. People who maintain fluoride use, low sugar frequency, and daily interdental cleaning keep more natural teeth into older age.

\section*{Prevention and Self-Care}
\begin{itemize}
    \item Brush morning and night; clean between teeth once daily every day, not only when food is stuck
    \item Choose fluoride toothpaste and drink fluoridated tap water where available
    \item Limit sugary and acidic drinks; prefer water and milk; use a straw for occasional soft drinks and do not sip all day
    \item Replace toothbrushes or powered heads every three months or when bristles splay
    \item Do not scrub horizontally with a hard brush; use gentle systematic coverage of all surfaces
    \item Clean around braces, retainers, and bridges with specialised tools as your dentist demonstrates
    \item Attend dental visits even if you feel fine; prevention is cheaper and kinder than emergency care
    \item Ask about access schemes if cost is a barrier---public dental services and child dental benefits may apply in many countries
\end{itemize}

\section*{Sources and Further Reading}
\begin{itemize}
    \item NHS. \textit{How to keep your teeth clean.} https://www.nhs.uk/
    \item Mayo Clinic. \textit{Oral health.} https://www.mayoclinic.org/
    \item MSD Manual. \textit{Dental care.} https://www.msdmanuals.com/
    \item MedlinePlus. \textit{Dental health.} https://medlineplus.gov/
    \item Better Health Channel. \textit{Teeth cleaning.} https://www.betterhealth.vic.gov.au/
    \item Healthdirect. \textit{Dental care.} https://www.healthdirect.gov.au/
\end{itemize}
